\documentclass[11pt]{article}
\usepackage[margin=1.05in]{geometry}
\usepackage{amsmath,amssymb}
\usepackage{booktabs}
\usepackage{graphicx}
\usepackage[T1]{fontenc}
\usepackage[utf8]{inputenc}
\usepackage[round,authoryear]{natbib}
\IfFileExists{lmodern.sty}{\usepackage{lmodern}}{}
\IfFileExists{microtype.sty}{\usepackage{microtype}}{}
\usepackage[hidelinks]{hyperref}
\usepackage{setspace}
\setstretch{1.12}
\newcommand{\E}{\mathbb{E}}

\title{The Macro Anchor as a Short-Rate Central Tendency:\\
Bond Risk Premia, Out-of-Sample Discipline, and a Specification Audit\thanks{%
Working paper. All results are produced by fully reproducible, deterministic notebooks
(expanding-window estimation throughout; FRED/OECD data with documented fallbacks). All
errors are mine. Comments welcome.}}
\author{Saverio Lauriola\thanks{Quantitative Risk Analyst. Email available on request. The
views expressed are the author's own and do not represent those of any employer.}}
\date{July 2026 --- first draft}

\begin{document}
\maketitle

\begin{abstract}
\noindent
This paper asks whether the intuition that nominal yields gravitate toward a
\emph{macro-consistent anchor} --- the neutral real rate plus trend inflation --- contains
tradeable information about bond risk premia, and answers it by escalating the discipline of
the test rather than the flexibility of the signal. Two natural formulations fail honestly:
long-yield levels barely mean-revert, and the raw 10-year ``anchor gap,'' despite strong
in-sample significance on a 50-year sample (Newey--West $t=2.86$) and cleanly monotone
return quintiles, does not survive out of sample ($R^2_{OOS}=-0.49$; Cochrane--Piazzesi
fares worse). The idea survives only in its theoretically disciplined form: the anchor as
the \emph{central tendency of a Hull--White short-rate process}, with reversion estimated at
the front end where it exists and the term premium obtained as an \emph{output}. The
resulting model term premium is not the slope in disguise (correlation $0.35$), tracks the
Kim--Wright benchmark ($0.70$) with three transparent parameters, is significant in sample
at every tenor with $t$-statistics monotone in maturity ($2.06$ to $3.87$), and is the only
predictor in a five-way horse race with a \emph{positive realized} out-of-sample $R^2$
(basket $+0.012$, Clark--West 10\%; exact 10Y target $+0.085$, Clark--West 5\%). A
pre-specified replication across nine developed markets yields $9/9$ positive slopes and
positive $R^2_{OOS}$ in $6/9$. The paper's distinctive section is a \emph{specification and
data audit} that decomposes the entire $0.012\!\to\!0.215$ spread across test variants into
named, quantified ingredients --- tenor concentration ($+0.07$), a duration-approximated
target ($+0.05$), the funding-rate series ($+0.07$), and sample length ($+0.06$) --- and
shows the reversion parameter is identified by pre-ZLB history. Results are robust to a
publication-lag inflation input. The framing throughout is \emph{theory as regularization}:
the signal with the best in-sample fit has the worst out-of-sample record, and the
model-disciplined form is the only survivor.
\end{abstract}

\vspace{0.5em}
\noindent\textbf{Keywords:} bond risk premia; term premium; affine short-rate models;
out-of-sample $R^2$; Clark--West; neutral rate; specification audit.

\section{Introduction}

A recurring intuition in macro-rates investing is that nominal yields should gravitate not
toward their own historical mean but toward a \emph{macro-consistent anchor}: the neutral
real short rate $r^*$ plus expected long-run inflation $\pi_e$ (plus, for long tenors, a
term premium). When the observed yield sits far above that anchor, bonds look cheap against
fundamentals; far below, rich. The intuition is old, respectable, and dangerous --- because
it invites exactly the kind of in-sample overfitting that the return-predictability
literature has learned to distrust \citep{goyalwelch,bauerhamilton}.

This paper takes the intuition seriously by doing the opposite of tuning it. We fix the
anchor's construction once, test the idea in three progressively more disciplined
formulations, report the failures with the same prominence as the success, and subject the
surviving result to a replication across nine markets and to a full specification and data
audit. Four findings emerge.

\textbf{First, two honest failures.} As a \emph{level-timing} signal the idea is dead on
arrival: at monthly frequency the 10-year level and slope are statistically
indistinguishable from random walks (estimated reversion speeds $\kappa\approx 0.02$ and
$\approx 0$). As a \emph{raw risk-premium} predictor --- regressing next-12-month bond
excess returns on the gap between the 10-year yield and the anchor, in the spirit of
\citet{cieslakpovala} --- the signal is strongly significant in sample on a five-decade
sample (Newey--West $t=2.86$, $R^2=0.13$, and mean subsequent excess returns rising
monotonically from $+0.13\%$ to $+6.55\%$ across gap quintiles) yet posts an out-of-sample
$R^2$ of $-0.49$ against the expanding-mean benchmark of \citet{campbellthompson}. For
context, the \citet{cochranepiazzesi} forward-rate factor does worse ($-1.24$) on the same
protocol. This is the \citet{goyalwelch} pattern in miniature, and we report it as a
result, not a nuisance.

\textbf{Second, the idea survives in exactly one form --- the theoretically disciplined
one.} We embed the anchor where the data say reversion actually lives: as the time-varying
central tendency $\theta_t=r^*_t+\pi_{e,t}$ of a Hull--White short-rate process
\citep{hullwhite}, estimated under the physical measure by regressing $\Delta r$ on the gap
to the anchor --- the project's original ``policy-gap'' intuition promoted to a parameter
estimator. The pure-expectations curve (expectations plus convexity, zero price of risk)
then delivers a \emph{model-implied term premium per tenor as an output},
$TP(\tau,t)=y_{obs}(\tau,t)-y_{PE}(\tau,t)$. Three honesty checks discipline the object:
the HW term premium correlates only $0.35$ with the slope (constant-$\theta$ Vasicek and
CIR analogues, lacking the anchor, correlate $0.77$ and $0.89$ --- they \emph{are} the
slope in disguise); it tracks the Kim--Wright five-factor benchmark at $0.70$ with three
transparent parameters (the anchor-free models track it at $0.03$ and $0.18$); and its
in-sample predictive $t$-statistics are monotone in tenor, from $2.06$ at two years to
$3.87$ at ten.

\textbf{Third, the surviving signal passes the out-of-sample gate that everything else
fails, and replicates across markets.} In a five-way horse race on a common evaluation
window (1997--2026, 336 months) the HW term premium is the only predictor with a positive
\emph{realized} out-of-sample $R^2$ ($+0.012$ on a 2--10Y basket; $+0.085$ on the exact
10-year target), with Clark--West significance at the 10\% and 5\% level respectively;
slope and Fama--Bliss have comparable Clark--West statistics but negative realized
$R^2_{OOS}$ --- population content destroyed by estimation noise --- and lose their 5\%
stars once evaluated on the common window. Running the \emph{same, pre-specified} pipeline
on nine developed markets yields positive slopes in $9/9$ countries, $R^2_{OOS}>0$ in
$6/9$, and Clark--West above the 10\% threshold in $7/9$, with the failures (Australia,
Norway, a too-short Swiss sample) reported alongside.

\textbf{Fourth --- and the paper's distinctive contribution --- an audit, not just a
result.} The panel's US row is far stronger ($R^2_{OOS}=+0.215$) than the strict US
benchmark ($+0.012$). Rather than headline the larger number, we decompose the entire gap
with a reconciliation ladder and a month-by-month data audit: tenor concentration accounts
for $+0.07$ (the effect lives at the long end; the basket dilutes it), the
duration-approximated target for $+0.05$ (it shares the carry with the predictor and omits
roll-down), the funding-rate series for $+0.07$ (T-bill versus interbank --- the raw
ten-year yields are \emph{numerically identical} across sources, so the residual is
entirely the short leg), and the longer sample for $+0.06$. The signal itself is identical
across specifications (the cross-variants coincide to machine precision). A sample-cut
diagnostic adds a structural finding: re-estimating on post-1997 data alone drives the
reversion speed to its floor ($\kappa\to 10^{-4}$) and the result collapses --- the anchor
parameter is \emph{identified by pre-ZLB history}; the zero-rate decade alone contains no
information about it. A pre-specified reading rule, written before the audit numbers, then fixes
the framing: the panel is long-sample cross-country replication; the strict benchmark is
the US common-window result.

The paper relates to the bond risk-premium literature descending from \citet{famabliss}
and \citet{cochranepiazzesi}, to the trend-inflation ``cycle'' factor of
\citet{cieslakpovala} --- our raw gap is essentially their object, and our failure/success
pattern refines where its content resides --- and to the out-of-sample skepticism of
\citet{goyalwelch}, \citet{bauerhamilton} and \citet{duffee}. Methodologically we lean on
\citet{campbellthompson} and \citet{clarkwest} for evaluation, on \citet{neweywest} for
inference with overlapping targets, on \citet{vasicek}, \citet{cir} and \citet{hullwhite}
for the model, on \citet{hlw} and \citet{kimwright} for the anchor's external references,
and on \citet{pesaransmith} for the panel mean-group estimate. Our contribution is not a
new predictor class; it is a demonstration, unusually complete in its accounting, of
\emph{where} a popular macro intuition holds, where it does not, and how much of any
headline number is specification.

\section{Data and the macro anchor}\label{sec:data}

\textbf{United States.} Monthly, June 1976 to May 2026 (599 observations). Constant-maturity
Treasury yields GS1--GS10 (FRED); the state variable is the 3-month T-bill (TB3MS); CPI is
CPIAUCSL; the external term-premium reference is the Kim--Wright 10-year series
(THREEFYTP10). The 1976 start is imposed by GS2 availability. All series are month-end.

\textbf{The anchor.} We deliberately freeze a simple, replicable construction:
\begin{equation}
\theta_t \;=\; r^*_t + \pi_{e,t},\qquad
\pi_{e,t}=\text{EWMA}_{36m}\!\left(\Delta_{12}\log \text{CPI}\right),\qquad
r^*_t=\text{EWMA}_{120m}\!\left(r_t-\pi_{e,t}\right),
\end{equation}
i.e.\ trend inflation is a three-year-half-life average of trailing year-on-year CPI
(the construction underlying the \citealp{cieslakpovala} cycle), and the neutral real rate
is a slow ten-year-half-life average of the realized real short rate (the
Holston--Laubach--Williams estimate is used instead whenever reachable; the notebook prints
which source was used). No term-premium block enters the \emph{short-rate} anchor: for long
tenors the term premium is an output of the model, not an input. Section~\ref{sec:robust}
shows the results are essentially unchanged when CPI enters with a one-month publication
lag.

\textbf{G10 panel.} Nine markets (USA, DEU, GBR, CAN, CHE, SWE, NOR, AUS, NZL; Japan is
dropped for insufficient overlapping data --- a convenient absence we flag, since Japan is
the archetypal ZLB case where the model is least at home). Sources are OECD MEI series on
FRED: 10-year government yields (\texttt{IRLTLT01*}), 3-month interbank rates
(\texttt{IR3TIB01*}), CPI (\texttt{*CPIALLMINMEI}, quarterly fallback for AUS/NZL). Spans
range from 1960 (DEU, GBR, CAN) to 1999 (CHE). Absent full foreign curves, the panel target
is the standard duration approximation
$rx^{(10)}_{t\to t+12}\approx y^{(10)}_t - D(y)\,\Delta y^{(10)} - r^{3m}_t$; the audit of
Section~\ref{sec:audit} quantifies exactly how much this construction flatters results.

\section{Two honest failures}\label{sec:failures}

\textbf{Level timing.} Estimated at monthly frequency, the reversion speeds of the first two
principal components of the US curve are $\kappa\approx 0.022$ (level) and $\approx 0$
(slope): half-lives of decades. Whatever the anchor knows, it is not that the ten-year
yield will fall soon. The allocation experiments that motivated this project (not reported
here; available in the replication repository) confirm that carry and roll-down, not level
reversion, do the work in curve strategies.

\textbf{The raw risk-premium gap.} Reframed as a return predictor, the natural but
ultimately failed version is the deviation of the 10-year yield from a legacy long-yield
anchor $r^*+\pi_e+TP$ regressed on next-12-month bond excess returns. This is not the
short-rate anchor used below: here the term-premium block is an external long-yield
benchmark component, whereas in the disciplined specification the term premium is derived
as a model output. On 1976--2026 the in-sample evidence is genuinely strong: Newey--West $t=2.86$ (lag 12), $R^2=0.13$, and mean
subsequent 12-month excess returns by gap quintile of
$\{+0.13,\,+1.9,\,+3.4,\,+4.7,\,+6.55\}$ percent --- clean monotonicity from ``rich'' to
``cheap.'' Out of sample, against the expanding-mean benchmark with expanding-window
re-estimation: $R^2_{OOS}=-0.49$. The Cochrane--Piazzesi forward factor, evaluated
identically, posts $-1.24$. On the modern subsample (2004--2026) the gap's in-sample
$t$-statistic is only $1.18$ with $R^2_{OOS}=-0.14$. We record both failures at full
prominence: the in-sample strength is concentrated in the 1980s disinflation and does not
transfer to a real-time forecaster. The natural response is not to tune the signal but to
ask \emph{where} the economics actually operate.

\section{The anchor as a short-rate central tendency}\label{sec:model}

Reversion is real at the front end: policy drags the short rate toward neutral. We
therefore place the anchor where the reversion lives, as the central tendency of the
instantaneous rate, and derive the whole curve.

\subsection{Models and estimation under the physical measure}

Three diffusions for the 3-month rate $r_t$:
Ornstein--Uhlenbeck/Vasicek $dr=\kappa(\theta-r)dt+\sigma dW$ (Gaussian transition, closed
form); CIR $dr=\kappa(\theta-r)dt+\sigma\sqrt{r}\,dW$ (non-central $\chi^2$ transition);
and Hull--White with the \emph{macro anchor as moving central tendency},
$dr=\kappa(\theta_t-r)dt+\sigma dW$ with $\theta_t=r^*_t+\pi_{e,t}$. Estimation is always
under the physical measure P from the time series --- never calibrated to today's curve, so
the comparison with the observed curve is informative rather than circular. OU is fit by
its exact AR(1) discretization; CIR by exact-moment weighted least squares, refined by
non-central-$\chi^2$ maximum likelihood; HW by regressing $\Delta r_t$ on the gap
$(\theta_t-r_t)$ --- the slope \emph{is} $\kappa\Delta$. All downstream signals use
\emph{expanding-window} parameters (120-month burn-in); full-sample estimates
(Table~\ref{tab:params}) are descriptive only.

\begin{table}[t]\centering
\caption{Full-sample physical-measure estimates, US 3-month rate, 1976--2026 (annualized).}
\label{tab:params}
\begin{tabular}{lcccc}
\toprule
Model & $\kappa$ & half-life (months) & $\theta$ & $\sigma$ \\
\midrule
OU / Vasicek & 0.094 & 88.8 & 3.87\% & 1.49\% \\
CIR (NC$\chi^2$ ML) & 0.053 & 157 & 3.57\% & 0.061 \\
HW (anchor) & 0.175 & 47.5 & $r^*_t+\pi_{e,t}$ (moving) & 1.48\% \\
\bottomrule
\end{tabular}
\par\smallskip
\footnotesize Notes: the CIR Feller condition is violated ($2\kappa\theta=0.0019<\sigma^2=0.0034$),
consistent with a decade at the zero bound. AR-based $\hat\kappa$ is biased upward near a
unit root. The HW gap-regression $t$-statistic is $1.37$: the halved half-life relative to
the constant-$\theta$ OU is a point estimate consistent with the thesis, not a
statistically decisive one --- we say so.
\end{table}

\subsection{The term premium as an output, and three honesty checks}

For each model the pure-expectations yield (expectations plus convexity, zero price of
risk) is closed-form affine, $y_{PE}(\tau)=\big(-\ln A(\tau)+B(\tau)\,r_t\big)/\tau$ with
$B(\tau)=(1-e^{-\kappa\tau})/\kappa$; the HW variant evaluates the Vasicek formula at
$\theta=\theta_t$ (a random-walk forecast of a slow anchor, stated rather than hidden). The
model term premium is the wedge $TP_m(\tau,t)=y_{obs}(\tau,t)-y^{m}_{PE}(\tau,t)$.
Table~\ref{tab:honesty} reports the checks that decide whether this object is new
information.

\begin{table}[t]\centering
\caption{Is the model term premium the slope in disguise? Does it track the benchmark?}
\label{tab:honesty}
\begin{tabular}{lcc}
\toprule
 & corr$(TP_{10},\ \text{slope})$ & corr$(TP_{10},\ \text{Kim--Wright})$ \\
\midrule
HW (anchor) & \textbf{0.35} & \textbf{0.70} \\
OU (constant $\theta$) & 0.77 & 0.03 \\
CIR & 0.89 & 0.18 \\
\bottomrule
\end{tabular}
\par\smallskip
\footnotesize Notes: the anchor is the only ingredient distinguishing HW from OU. Without
it, the ``term premium'' is a repackaged slope and tracks nothing; with it, a
three-parameter transparent model tracks the Fed's five-factor benchmark at $0.70$.
\end{table}

\section{Predictive evidence}\label{sec:pred}

Targets are per-tenor 12-month excess returns in the Fama--Bliss construction,
$rx^{(n)}_{t\to t+12}=n\,y_t(n)-(n\!-\!1)\,y_{t+12}(n\!-\!1)-y_t(1)$ (off-grid maturities
interpolated), and a 2--10Y equal-weight basket. Inference uses Newey--West with lag 12
(overlapping targets). Out-of-sample $R^2$ follows \citet{campbellthompson}: expanding
forecasting regressions against the expanding mean. Nested significance uses
\citet{clarkwest} with NW(12) errors, one-sided.

\begin{table}[t]\centering
\caption{In-sample predictive regressions: $rx^{(n)}\sim TP_{HW}(n)$, Newey--West lag 12.}
\label{tab:tenor}
\begin{tabular}{lccccc}
\toprule
Tenor & 2y & 3y & 5y & 7y & 10y \\
\midrule
$\beta$ per $+1$SD (\%/12m) & $+0.25$ & $+0.48$ & $+1.24$ & $+2.08$ & $+3.48$ \\
NW $t$ & 2.06 & 2.10 & 2.80 & 3.35 & \textbf{3.87} \\
$R^2$ & 0.035 & 0.037 & 0.076 & 0.117 & 0.172 \\
\bottomrule
\end{tabular}
\par\smallskip
\footnotesize Notes: basket $t=3.04$, $R^2=0.087$. The monotone tenor structure is
documented here, \emph{before} any out-of-sample claim; it licenses reporting the exact
10Y result in Table~\ref{tab:horse} as structure rather than selection.
\end{table}

\begin{table}[t]\centering
\caption{Horse race on the common evaluation window (1997:05--2026:05, $n=336$ for all).}
\label{tab:horse}
\begin{tabular}{lccc}
\toprule
Predictor & $R^2_{OOS}$ & Clark--West $t$ & significance \\
\midrule
$TP_{HW}$ (anchor model), basket & $\;\;\,+0.012$ & 1.31 & 10\% \\
$TP_{HW}$, exact 10Y target & $\;\;\,+0.085$ & 2.25 & \textbf{5\%} \\
$TP_{OU}$ (constant $\theta$) & $-0.080$ & 1.19 & --- \\
Slope $(y_{10}-y_1)$ & $-0.010$ & 1.60 & 10\% \\
Fama--Bliss forward spread & $-0.015$ & 1.62 & 10\% \\
Raw anchor gap (legacy) & $-0.468$ & 0.51 & --- \\
\bottomrule
\end{tabular}
\par\smallskip
\footnotesize Notes: on their longer native windows, slope and Fama--Bliss carry 5\%
Clark--West stars ($t=1.66,\,1.74$) with \emph{negative} realized $R^2_{OOS}$ ---
population-level content eaten by estimation noise; both stars fail to survive the common
window. The HW term premium is the only predictor whose realized forecasts beat the
benchmark. Adding the slope to $TP_{HW}$ \emph{worsens} $R^2_{OOS}$ ($-0.066$): the
information is in the model object.
\end{table}

The claims hierarchy we defend is therefore: (i) the pre-specified strict benchmark ---
basket $R^2_{OOS}=+0.012$, Clark--West 10\%; (ii) the documented tenor structure --- exact
10Y $R^2_{OOS}=+0.085$, Clark--West 5\%; and, next, (iii) cross-country replication with
magnitudes explicitly discounted by the audit of Section~\ref{sec:audit}.

\section{Cross-country replication}\label{sec:panel}

The identical pipeline --- same anchor recipe and half-lives, same HW gap regression, same
expanding windows, same tests, nothing re-tuned --- is run on nine developed markets.
Because nothing is chosen on the new data, this stage is better interpreted as a frozen-pipeline
replication exercise than as another draw from a specification search. Cross-country dependence
and heterogeneous sample spans still warrant caution, so the sign-replication count is treated
as the primary panel statistic.

\begin{table}[t]\centering
\caption{G10 replication of the pre-specified pipeline (duration-approximated 10Y target).}
\label{tab:panel}
\begin{tabular}{lcccccc}
\toprule
 & $\beta$/SD & NW $t$ & $R^2_{IS}$ & $R^2_{OOS}$ & CW $t$ & $n_{OOS}$ \\
\midrule
USA & 0.047 & 5.64 & 0.29 & $+0.216$ & 4.74 & 467 \\
CAN & 0.036 & 5.18 & 0.20 & $+0.222$ & 4.69 & 520 \\
SWE & 0.032 & 4.27 & 0.21 & $+0.201$ & 3.20 & 197 \\
DEU & 0.025 & 3.42 & 0.11 & $+0.016$ & 1.73 & 520 \\
NOR & 0.024 & 3.18 & 0.14 & $-0.080$ & 1.65 & 221 \\
AUS & 0.017 & 2.00 & 0.04 & $-0.112$ & $-0.49$ & 404 \\
NZL & 0.027 & 1.90 & 0.07 & $+0.110$ & 3.06 & 350 \\
CHE & 0.019 & 1.78 & 0.11 & $-0.429$ & 0.38 & 47 \\
GBR & 0.014 & 1.50 & 0.03 & $+0.023$ & 1.79 & 520 \\
\midrule
\multicolumn{7}{l}{\footnotesize Signs: \textbf{9/9} positive. NW $t>2$: 5/9 ($>1.65$: 8/9).
$R^2_{OOS}>0$: \textbf{6/9}. CW $>10\%$ crit.: 7/9.}\\
\multicolumn{7}{l}{\footnotesize Mean-group $\beta=+2.66\%$/SD, $t=7.85$
\citep{pesaransmith} --- overstated by cross-country dependence;}\\
\multicolumn{7}{l}{\footnotesize the sign-replication count is the primary statistic, as
pre-specified. Failures reported: AUS, NOR; CHE has}\\
\multicolumn{7}{l}{\footnotesize only 47 evaluations; JPN dropped for data overlap (the ZLB
case --- its absence is flagged, not hidden).}\\
\bottomrule
\end{tabular}
\end{table}

\section{The specification and data audit}\label{sec:audit}

The panel's US row ($R^2_{OOS}=+0.215$) is an order of magnitude stronger than the strict
US benchmark ($+0.012$). A result this sensitive to specification is not reportable until
the sensitivity itself is measured. We do so in three steps, all on identical evaluation
months, with reading rules written \emph{before} the numbers.

\subsection{A reconciliation ladder: target versus signal}

On the same US data and common window we evaluate five pairs: (A1) basket signal
$\to$ basket exact target; (A2) 10Y signal $\to$ exact 10Y target; (B1) panel-style signal
$\to$ duration-approximated 10Y target; and the cross-variants (B2) affine signal on the
approximated target, (B3) panel signal on the exact target.
Table~\ref{tab:ladder} shows B1$\equiv$B2 and B3$\equiv$A2 \emph{to machine precision}: the
panel's signal is numerically the same object as the affine notebook's --- so the entire
on-sample gap is the \textbf{target}, never the signal. The decomposition is two steps:
basket $\to$ exact 10Y adds $+0.07$ of $R^2_{OOS}$ (tenor concentration, licensed by
Table~\ref{tab:tenor}); exact $\to$ duration-approximated adds $+0.05$. The approximation
flatters for two identifiable reasons: it shares the yield level (carry) mechanically with
the predictor, and, lacking roll-down, aging and cross-tenor noise, it is a smoother target
that a slow signal explains more easily. It is the standard construction of the
international bond-premium literature, but the honest hierarchy is exact-Fama--Bliss
first.

\begin{table}[t]\centering
\caption{Reconciliation ladder, US data, common window (1997:05--, $n=336$).}
\label{tab:ladder}
\begin{tabular}{lccc}
\toprule
Pair & $R^2_{OOS}$ & CW $t$ & NW $t$ \\
\midrule
A1: basket signal $\to$ basket exact & $+0.012$ & 1.31 & 3.04 \\
A2: 10Y signal $\to$ 10Y exact & $+0.085$ & 2.25 & 3.87 \\
B1: panel signal $\to$ 10Y dur.-approx. & $+0.140$ & 2.70 & 4.31 \\
B2: 10Y signal $\to$ 10Y dur.-approx. & $+0.140$ & 2.70 & 4.31 \\
B3: panel signal $\to$ 10Y exact & $+0.085$ & 2.25 & 3.87 \\
\bottomrule
\end{tabular}
\par\smallskip
\footnotesize Notes: B1$=$B2 and B3$=$A2 exactly --- the signal is invariant; the target
explains everything on-sample.
\end{table}

\subsection{A row-by-row data audit: the funding series}

What remains between the affine-data B1 ($+0.135$ on common rows) and the panel-USA row is
data source and sample. A month-by-month comparison of every ingredient over the common
1976--2025 span (586 months) finds: the ten-year yields are \emph{numerically identical}
(the OECD long rate is the Fed's GS10; correlation $1.000$, zero difference); CPI is
identical to a few basis points. The entire source residual is the \textbf{3-month rate}:
T-bill versus interbank, mean difference $-53$bp, spiking to $-365$bp in October 2008 (the
TED blowout) and to $-250/-316$bp in 1980--82 --- and it propagates into the $r^*$ proxy
(correlation drops to $0.70$, also via EWMA burn-in state), the anchor ($-94$bp mean), the
model TP ($+67$bp) and the target's funding leg. On \emph{identical rows} the same formula
yields $R^2_{OOS}=0.135$ (CW $2.67$) with the T-bill and $0.208$ (CW $4.17$) with the
interbank rate: \textbf{seven points of out-of-sample $R^2$ hinge on the choice of the
short-rate series}. We report both and pick neither; there is a defensible case for the
interbank rate as the state variable (a smoother funding rate, free of T-bill
flight-to-quality distortions), and an equally standard case for the T-bill as the correct
funding leg for Treasury excess returns. The final ingredient, extending the sample from
1976 back to 1964, is worth roughly $+0.06$.

\subsection{Sample cuts and parameter identification}

Re-running the panel pipeline on US windows: from 1976 the result survives
($R^2_{OOS}=+0.153$, CW $3.03$) --- it is not a pre-1976 artifact. Re-estimating on
post-1997 data \emph{alone}, it collapses ($R^2_{OOS}=-0.061$, $n=72$) with the reversion
speed at its numerical floor ($\hat\kappa=10^{-4}$). The distinction with the strict
benchmark matters: the benchmark \emph{learns} $\kappa$ from 1976 onward and
\emph{evaluates} from 1997 (and survives); the cut both learns and evaluates post-1997.
The economics: \textbf{the anchor-reversion parameter is identified by the pre-ZLB
history} --- a decade pinned at zero contains no information about reversion toward a
moving neutral rate. The pre-specified rule then fixes the framing: the panel is
long-sample cross-country replication; the US common-window result is the strict
benchmark.

\subsection{Accounting for the full gap}

Adding the pieces: tenor concentration $+0.07$, target approximation $+0.05$, funding
series $+0.07$, sample $+0.06$ --- the entire spread from the strict $+0.012$ to the panel
US $+0.215$ is attributed to named, quantified specification choices, with the signal
itself invariant. To our knowledge such an accounting is rarely reported in this
literature; we suggest it should be standard whenever a predictability claim spans several
test designs.

\section{Robustness and the distributional layer}\label{sec:robust}

\textbf{Publication-lag inflation.} Shifting CPI by one month (the real-time proxy) leaves
the strict benchmark essentially unchanged: $R^2_{OOS}$ $+0.0093$ versus $+0.0116$,
Clark--West $1.32$ versus $1.31$, NW $t$ $3.07$. The 36-month-half-life trend makes the
signal slow by construction; this verifies it. Full real-time treatment (ALFRED vintages,
revisions and not just lag) is the stated upgrade, not claimed.

\textbf{Calibration of the transition densities.} Probability-integral-transform tests of
the 12-month-ahead short-rate distribution (expanding parameters) favor the
level-dependent-volatility CIR over the Gaussian OU (KS distance $0.093$ versus $0.209$;
overlapping observations, so shapes rather than stars). An honest negative rides along:
in a variance-forecast race, trailing realized variance dominates both models'
closed-form variances ($R^2_{OOS}$ $0.685$ versus $0.013$ for CIR and $-0.087$ for OU) ---
the ``closed forms win on second moments'' conjecture fails against simple persistence.
The two models' 1\% lower quantiles diverge sharply as rates approach zero (maximum
disagreement September 2011), a clean model-based ZLB regime flag.

\textbf{Economic value.} Sizing duration on the model's closed-form $\E[rx]$ and $SD[rx]$
(constant-TP forecast for the aged yield; signals lagged one month) improves on static
duration: Sharpe $0.58\to 0.71$--$0.75$ for sign/linear sizing with maximum drawdown
halved ($-19.7\%\to -8$/$-11\%$); the $z$-sized book is the most defensive
(vol $2.3\%$, MaxDD $-6.1\%$, Sharpe $0.63$). Gross of transaction costs.

\section{Discussion: theory as regularization, and multiple testing}\label{sec:disc}

Two patterns organize the results. First, \emph{theory as regularization}: the raw anchor
gap has the best in-sample fit ($R^2=0.128$) and the worst out-of-sample record
($-0.47$); the model-disciplined form has less in-sample fit and is the only out-of-sample
survivor. Imposing the process structure --- expectations, convexity, reversion placed
where it is identified --- disciplines the signal the way a prior disciplines a regression.
Second, an explicit multiple-testing account: this is the \emph{third} formulation of one
idea tested on US data (the two failures are reported at equal prominence), so the
marginal significance levels deserve a further mental discount \citep{harveyliuzhu}; the
G10 stage, being a frozen pipeline on new data, does not add to the count. We regard the
resulting claim --- ``survives out of sample in its disciplined form, with a marginal star,
replicated in sign across nine markets, with every specification sensitivity quantified''
--- as the strongest \emph{honest} statement the data permit, and deliberately weaker than
the strongest statement the numbers would superficially allow.

\section{Limitations}\label{sec:limits}

One-factor models cannot fit a three-factor curve; part of any model TP is misfit and is
treated relative to its own history. $\hat\kappa$ is biased upward near a unit root, and
its gap-regression $t$-statistic (1.37) is suggestive, not decisive. The CIR Feller
condition is violated in a ZLB sample, and the Gaussian HW admits negative rates. The
duration-approximated panel target flatters ($+0.05$, quantified); panel data are revised
OECD series with heterogeneous spans; cross-country observations are not independent.
Results are sensitive to the funding-rate series ($\pm 0.07$; both reported). Economic
value is gross of costs. The evaluation windows, while common where it matters, remain a
single historical path.

\section{Conclusion}\label{sec:concl}

A popular macro intuition --- yields gravitate toward $r^*$ plus trend inflation ---
contains real information about bond risk premia, but only when embedded where its
economics operate: as the central tendency of the short rate, with the term premium
derived rather than assumed. Tested with pre-specified protocols, the disciplined form is
the only predictor among standard benchmarks whose realized out-of-sample forecasts beat
the expanding mean, replicates in sign across nine markets, survives a real-time inflation
proxy, and comes with a complete specification audit that attributes every point of
variation across test designs. The failures along the way --- level timing, the raw gap,
the variance-forecast conjecture, three countries out of sample, a decade that identifies
nothing --- are reported at the same volume as the successes, because in this literature
the reporting standard \emph{is} the result.

\begin{thebibliography}{99}\small

\bibitem[Adrian et al.(2013)]{acm} Adrian, T., R.~K. Crump, and E.~Moench (2013).
Pricing the term structure with linear regressions. \emph{Journal of Financial
Economics} 110(1), 110--138.

\bibitem[Bauer and Hamilton(2018)]{bauerhamilton} Bauer, M.~D., and J.~D. Hamilton (2018).
Robust bond risk premia. \emph{Review of Financial Studies} 31(2), 399--448.

\bibitem[Campbell and Thompson(2008)]{campbellthompson} Campbell, J.~Y., and S.~B.
Thompson (2008). Predicting excess stock returns out of sample: Can anything beat the
historical average? \emph{Review of Financial Studies} 21(4), 1509--1531.

\bibitem[Cieslak and Povala(2015)]{cieslakpovala} Cieslak, A., and P.~Povala (2015).
Expected returns in Treasury bonds. \emph{Review of Financial Studies} 28(10), 2859--2901.

\bibitem[Clark and West(2007)]{clarkwest} Clark, T.~E., and K.~D. West (2007).
Approximately normal tests for equal predictive accuracy in nested models.
\emph{Journal of Econometrics} 138(1), 291--311.

\bibitem[Cochrane and Piazzesi(2005)]{cochranepiazzesi} Cochrane, J.~H., and M.~Piazzesi
(2005). Bond risk premia. \emph{American Economic Review} 95(1), 138--160.

\bibitem[Cox et al.(1985)]{cir} Cox, J.~C., J.~E. Ingersoll, and S.~A. Ross (1985).
A theory of the term structure of interest rates. \emph{Econometrica} 53(2), 385--407.

\bibitem[Duffee(2002)]{duffee} Duffee, G.~R. (2002). Term premia and interest rate
forecasts in affine models. \emph{Journal of Finance} 57(1), 405--443.

\bibitem[Fama and Bliss(1987)]{famabliss} Fama, E.~F., and R.~R. Bliss (1987).
The information in long-maturity forward rates. \emph{American Economic Review} 77(4),
680--692.

\bibitem[Goyal and Welch(2008)]{goyalwelch} Goyal, A., and I.~Welch (2008).
A comprehensive look at the empirical performance of equity premium prediction.
\emph{Review of Financial Studies} 21(4), 1455--1508.

\bibitem[Harvey et al.(2016)]{harveyliuzhu} Harvey, C.~R., Y.~Liu, and H.~Zhu (2016).
\ldots and the cross-section of expected returns. \emph{Review of Financial Studies}
29(1), 5--68.

\bibitem[Holston et al.(2017)]{hlw} Holston, K., T.~Laubach, and J.~C. Williams (2017).
Measuring the natural rate of interest: International trends and determinants.
\emph{Journal of International Economics} 108, S59--S75.

\bibitem[Hull and White(1990)]{hullwhite} Hull, J., and A.~White (1990).
Pricing interest-rate-derivative securities. \emph{Review of Financial Studies} 3(4),
573--592.

\bibitem[Kim and Wright(2005)]{kimwright} Kim, D.~H., and J.~H. Wright (2005).
An arbitrage-free three-factor term structure model and the recent behavior of long-term
yields and distant-horizon forward rates. FEDS Working Paper 2005-33.

\bibitem[Kupiec(1995)]{kupiec} Kupiec, P.~H. (1995). Techniques for verifying the accuracy
of risk measurement models. \emph{Journal of Derivatives} 3(2), 73--84.

\bibitem[Newey and West(1987)]{neweywest} Newey, W.~K., and K.~D. West (1987).
A simple, positive semi-definite, heteroskedasticity and autocorrelation consistent
covariance matrix. \emph{Econometrica} 55(3), 703--708.

\bibitem[Pesaran and Smith(1995)]{pesaransmith} Pesaran, M.~H., and R.~Smith (1995).
Estimating long-run relationships from dynamic heterogeneous panels.
\emph{Journal of Econometrics} 68(1), 79--113.

\bibitem[Vasicek(1977)]{vasicek} Vasicek, O. (1977). An equilibrium characterization of
the term structure. \emph{Journal of Financial Economics} 5(2), 177--188.

\end{thebibliography}

\end{document}
